\chapter{Conclusion}
\label{chapter:Conclusion}

\section{Reflection on Results}

In this work, we studied whether the confidence produced by a decision-tree
signal filter was related to the out-of-sample performance of a rule-based
trading strategy in the gold market. More specifically, we considered the
question of how decision-tree confidence relates to the out-of-sample
performance of a multi-indicator strategy built on the broker symbol
\texttt{XAUUSD} on the H1 timeframe.

To answer this question, we implemented a reproducible Python-based experiment
pipeline that retrieved MT5 market data, constructed technical indicators,
generated candidate trade signals, trained a decision tree on labelled signal
outcomes, and evaluated both baseline strategies and confidence-filtered
variants on a held-out market segment. The final experiment used a three-year
H1 window from January 2022 to December 2024, with the out-of-sample
evaluation period beginning on 28 May 2024.

The results showed that the unfiltered multi-indicator strategy remained
profitable out of sample, but that the decision-tree confidence score provided
additional information about signal quality. The clearest result was obtained
for confidence Range~4, which outperformed the unfiltered multi-indicator
baseline on total return, Sharpe ratio, win rate, and profit factor. At the
same time, the confidence ranges did not form a strictly increasing ordering:
some higher-confidence groups performed weakly, including Range~5. The most
important conclusion from the study is therefore that decision-tree confidence
was useful for separating stronger from weaker candidate signals in the chosen
\texttt{XAUUSD} H1 setting, but that the relationship was informative rather
than perfectly monotonic.

This interpretation is broadly consistent with prior literature that motivates
simple and interpretable filtering layers as a way to improve rule-based
strategies without immediately moving to highly complex models. Our results do
not justify the stronger claim that a decision tree universally improves
trading performance. They do, however, provide evidence that an interpretable
classifier can contribute practical ranking information when applied to
candidate signals generated by a transparent rule-based strategy.

\section{Limitations}

This study has several limitations that should be considered when interpreting
its conclusions. First, the analysis is limited to the broker-specific
instrument \texttt{XAUUSD}, used in this thesis as the tradable representation
of the gold market, which means that the findings cannot automatically be
generalised to other markets, brokers, or instrument definitions. Second, the
study examines only data from the H1 timeframe and only one final three-year
window, so the results may differ across shorter and longer trading horizons
and across different market periods.

Third, the evaluation is conducted in a controlled out-of-sample backtesting
environment built in Python on MT5-sourced data. This improves reproducibility
and supports detailed inspection of the machine-learning workflow, but it does
not fully capture live-trading conditions such as execution uncertainty, broker
latency, or slippage beyond the fixed transaction-cost assumptions used in the
study. Fourth, the machine-learning component is deliberately restricted to a
simple decision tree in order to preserve interpretability and reduce model
complexity, meaning that the conclusions should primarily be understood within
this specific experimental setting rather than as a general statement about
machine learning in trading.

\section{Ethical Implications}

The conclusions of this study have ethical and societal implications because
they concern the use of machine learning in financial decision-making. Even
though the model used in this thesis is intentionally simple and interpretable,
the results could still contribute to increased confidence in automated trading
systems. This creates a risk that such systems are viewed as more reliable,
objective, or robust than they actually are, especially when conclusions are
based on historical out-of-sample backtesting rather than live trading.

In a broader societal context, increased reliance on automated decision support
in finance may reduce human oversight and encourage decision-making based on
models that remain uncertain and sensitive to changing market conditions. For
this reason, the results should be interpreted carefully and understood as a
research contribution on interpretable signal filtering rather than as financial
advice or proof of guaranteed future profitability.

\section{Considerations for Future Work}

Future work could examine whether the observed relationship between
decision-tree confidence and trading performance remains stable beyond the
specific setting used in this thesis. Since the analysis is limited to
\texttt{XAUUSD} on the H1 timeframe, an important next step would be to test
the same approach on other instruments, brokers, and trading horizons. This
would help determine whether the conclusions reflect a broader property of
confidence-based filtering or are mainly tied to the present experimental
setting.

Another relevant direction would be to assess robustness under changing market
conditions. Although this thesis uses an out-of-sample design, future studies
could apply more extensive walk-forward testing, rolling retraining, or regime-
based evaluation in order to assess whether the usefulness of the confidence
estimates remains stable over time. In addition, the exported MQL5
implementation could be validated more extensively in the MT5 Strategy Tester
and compared directly with the Python backtest outputs to study implementation
parity between research and deployment environments.

It would also be valuable to explore whether the results depend on the specific
model and trade-management design choices. The present study deliberately uses
a simple decision tree to preserve interpretability, but future research could
compare this approach with other interpretable models or with alternative
definitions of trade outcomes and exit rules. Such work would clarify whether
the findings of this thesis are primarily linked to the present design choices
or indicate a more general role for interpretable machine learning in
rule-based trading.

\section{Usage of AI Tools}

Large language models were used as support for language revision, code
refactoring assistance, and drafting help during the development of the
software artefact. All experimental design decisions, code execution, result
verification, and final thesis formulations were manually reviewed by the
authors.
